\section{Related work}

\paragraph{Linear relational embeddings.} Treating relations as linear maps
between concept representations goes back to
\citet{DBLP:conf/icml/PaccanaroH00};
\citet{hernandez2024linearity} carried the idea over
to transformer language models \citep{vaswani2017attention}, and
subsequent work has asked why the resulting LREs are faithful for some
relations and not others.
\citet{merullo2025linear} trace part of this
variation to pretraining data: LREs are more faithful when subjects and
objects frequently co-occur in the corpus (we compare against such
frequency measures in Appendix~\ref{sec:corpus-freq-baselines}). A further question is whether such
structure reflects what a model has learned or only the particular
parametrization that realizes it, since many parametrizations induce the
same next-token distribution. Existing results point to the former,
since under suitable conditions, relational linearity holds either in
all distribution-equivalent models or in none, and it has been shown
empirically, through a probe based on the Kullback--Leibler divergence,
to vary across models, layers and phrasings of the relation
\citep{marconato2024all, valer2026relational}.
These works evaluate and explain linearity at the level of the relation.
We ask instead how much of it can be predicted from how the model
represents the subject, using a readout estimated without the relation
or any of its objects.

\paragraph{Subject representations in factual recall.} Mechanistic studies
converge on the view that much factual knowledge is attached to the
subject before the relation is processed. Feed-forward layers operate as
key--value memories whose keys respond to patterns in the input and
whose values induce distributions over the output vocabulary, and in
middle layers at the last subject token they are where causal tracing
localizes factual associations and where targeted edits change
individual facts
\citep{geva2021transformer, meng2022locating}.
Building on these findings, factual recall has been decomposed into
stages: MLP layers enrich the subject representation with many of its
attributes, a lookup that depends largely on the subject's name rather
than on the preceding context; the relation is propagated to the final
position; and attention heads extract the requested attribute
\citep{geva-etal-2023-dissecting, nanda2024fact}.
\citet{10.1162/TACL.a.746} show that knowledge
acquisition tracks how often subject and answer co-occur, though
frequency alone does not fully explain learning, which is why we treat
the corpus-frequency baselines in Appendix~\ref{sec:corpus-freq-baselines} as a comparison rather
than an explanation. Our approach can be read as a probe of this
enrichment, since it reports what the model is prepared to say about a
subject before any relation is given. This allows us to connect two
lines of work: We hypothesize that the enrichment the lens diagnoses
contributes to the relational linearity found by
\citet{hernandez2024linearity}. The lens predicts LRE
faithfulness across relations in all seven models we test, and the
association persists when controlling for object cardinality and
answer-token frequency (§\ref{sec:results}).

\paragraph{How facts are stored.} Theoretical work offers two accounts of
how transformers store facts.
\citet{nichani2025understanding} model attention and MLP
weights as associative memories and prove that a one-layer transformer
can store a number of facts that scales linearly with its parameter
count, with storage traded off between self-attention and MLP
parameters. \citet{DBLP:journals/corr/abs-2605-12426} study a
geometric form of storage: in a controlled one-layer setting, subject
embeddings encode linear superpositions of their associated attributes,
and a small MLP selects among them conditioned on the relation. In
pretrained models the two need not compete: lookup-like computation can
enrich the subject representation over the early and middle layers, and
the geometric picture then describes the enriched state from which a
relation selects. Our hypothesis concerns this enriched state: a
relation should be linearly decodable to the extent that its attribute
is prominently encoded in the subject representation after enrichment.

\paragraph{Lenses.} Lenses read out intermediate hidden states in terms of
the output vocabulary and differ mainly in the map they apply: the
model's own unembedding
\citep{nostalgebraist2020interpreting}, an affine translator
learned per layer \citep{belrose2023eliciting}, or
the model's input--output Jacobian averaged over a large corpus and
applied to raw activations, with the aim of identifying representations
the model is poised to verbalize
\citep{gurnee2026workspace}.
Our entity lens instead evaluates the Jacobian at reference entities of
the subject's type under a single neutral factual prompt, so that the
readout concerns what the model is prepared to say about entities of
that type rather than about any concept. Its contrastive variant further
applies the Jacobian to differences between entity representations,
subtracting what the references share.



%\citet{DBLP:conf/icml/PaccanaroH00} introduce linear
%relational embeddings (LREs)
%and \citet{hernandez2024linearity} show that LREs are part
%of the mechanism by which transformer models answer factual questions.
%
%
%\citet{merullo2025linear} show that the frequency
%of subject-object cooccurrence
%in the
%pretraining corpus is an important determinant of
%linearity. In our work, we put forward the hypothesis that
%relational linearity is primarily a property of how the model
%represents the subject, independent of the relation (or object). We also
%identify object cardinality as an additional important
%variable for linearity. Finally, we analyze limitations of
%LRE for understanding linearity and make suggestions for how
%future factual recall datasets should be constructed for
%relational linearity research.
%
%\citet{geva-etal-2023-dissecting} proposes subject
%enrichment as a key mechanism for factual recall. We can
%interpret our contrastive Jacobian lens as a probe of
%subject enrichment.
%Crucially, we connect in this paper two different lines of
%research on the representation of relations in transformer
%models: we hypothesize that the type of subject enrichment that the Jacobian
%lens diagnoses gives rise to the linearity of relations
%found by \citet{hernandez2024linearity}.
%
%DISCUSS \cite{10.1162/TACL.a.746}
%
%DISCUSS \cite{DBLP:journals/corr/abs-2605-12426}

